\documentclass[12pt]{article}
\usepackage[utf8]{inputenc}
\usepackage[T1]{fontenc}
\usepackage{hyperref}
\title{Bambam Language Resources for K-12 Curriculum Development}
\author{Compiled by Claude Sonnet 5.5 (claude-sonnet-5-5) for Andr\'as Kornai}
\date{}
\begin{document}
\maketitle

\section{Introduction}
Bambam (ISO 639-3 ptu; Glottocode bamb1270), also called Pitu Ulunna Salu, is an Austronesian language of West Sulawesi, Indonesia, in the Mamasa area. Speaker figures vary between sources. Searches found a few descriptive and reference pages but no school materials.

\section{Primary Educational Resources}
No primers, textbooks or school programmes were found in the search results.

\section{Linguistic and Reference Materials}
\begin{itemize}
\item Wikipedia, ``Bambam language'' (overview, phonology, and references to P.\ J.\ Campbell's 1989 MA thesis on Pitu Ulunna Salu grammar and his 1991 phonology paper; the full texts were not found online): \url{https://en.wikipedia.org/wiki/Bambam_language}
\item Grambank entry for Bambam: \url{https://grambank.clld.org/languages/bamb1270}
\item Omniglot, Bambam (writing system and notes): \url{https://omniglot.com/writing/bambam.htm}
\item Wikipedia, ``Mamasa language'' (a neighbouring language, for context): \url{https://en.wikipedia.org/wiki/Mamasa_language}
\end{itemize}

\section{Religious and Community Resources}
Omniglot mentions a New Testament translation (2004); I did not verify its details. Global Recordings Network has a language page:
\begin{itemize}
\item \url{https://www.globalrecordings.net/en/language/ptu}
\end{itemize}

\section{Digital and Audio Resources}
The Global Recordings Network page above is the only audio-related lead found. No dictionary or corpus was found.

\section{Curriculum Development Context}
No evidence of use in schools was found. Campbell's grammar and phonology work would be the main starting points for orthography and grammar, but they need to be located through libraries. Any curriculum would require new community-led materials.
\end{document}